\documentclass[11pt,a4paper]{article}
\usepackage[margin=2.5cm]{geometry}
\usepackage{longtable,array,booktabs,enumitem,hyperref}
\setlength{\parindent}{0pt}
\setlength{\parskip}{6pt}

\title{VeriPilot: Decisions for the Whole Team}
\author{Team of five $\cdot$ with the industrial mentor}
\date{Draft started 30 Oct 2026}

\begin{document}
\maketitle

Each item below needs all five of us to agree. Fill in the \textbf{Decision}, \textbf{Date} and \textbf{Owner} once settled. Items marked \emph{Blocking} must be closed before the stated milestone.

\section*{How we decide}
\begin{itemize}[nosep]
  \item Propose an option, discuss for a few days, then decide together. Record the result here.
  \item Anything about licensing, naming or making the repo public is decided by all five, not by one person.
  \item Ask the industrial mentor for input on items marked (M).
\end{itemize}

\section{Blocking before the repository goes public}
\begin{longtable}{p{3.2cm} p{5.3cm} p{2.2cm} p{3.8cm}}
\toprule
\textbf{Decision} & \textbf{Options / notes} & \textbf{Blocking?} & \textbf{Result (date, owner)} \\
\midrule
\endhead
Licence for our own code & MIT (simplest, permissive) or Apache-2.0 (adds a patent grant). Must be compatible with the licences of tools we reuse. Check each repo's licence first; do not copy code from unlicensed repos. & Before public push & \\
Project name & Keep ``VeriPilot'' only if GitHub and PyPI searches show it is free. Fallbacks in the proposal: BugBloom, CoverQuest, SeedSense. & Before public push & \\
Where the repo lives & Personal account of one member, or a GitHub organisation (better for five owners). & Before public push & \\
Contribution rules & Pull requests only, one reviewer, protected \texttt{main}. Do we require a sign-off (DCO) on commits? & Before first PR & \\
\bottomrule
\end{longtable}

\section{Technical decisions}
\begin{longtable}{p{3.2cm} p{5.3cm} p{2.2cm} p{3.8cm}}
\toprule
\textbf{Decision} & \textbf{Options / notes} & \textbf{Needed by} & \textbf{Result (date, owner)} \\
\midrule
\endhead
Benchmark designs (M) & 4 to 6 small open designs, e.g.\ UART, SPI, FIFO, timer. Must run on Verilator, so avoid heavy UVM. & Week 4 (13--20 Nov) & \\
Planted-bug protocol (M) & How many bugs (10--20 per design), which types, who plants them, how we record the ground truth. & Week 3 & \\
Linters to support & Verilator \texttt{--lint-only} first. Add a second linter (e.g.\ Verible) later? & Week 3 & \\
Python version and tooling & Python 3.10 on current machines. Virtual environment, dependency file, formatter and linter for our own Python. & Week 2 & \\
Testbench style & Plain SystemVerilog, cocotb (Python), or both. & Week 2--3 & \\
Language model for fixes & Small open model that runs locally so design code never leaves the machine. Which model and how it is run. & Build phase & \\
Baseline protocol & Exactly how setup A (normal way) is measured, so B and C are fair comparisons. Numeric targets only after the first baseline run. & Week 3--4 & \\
\bottomrule
\end{longtable}

\section{Team and process}
\begin{longtable}{p{3.2cm} p{5.3cm} p{2.2cm} p{3.8cm}}
\toprule
\textbf{Decision} & \textbf{Options / notes} & \textbf{Needed by} & \textbf{Result (date, owner)} \\
\midrule
\endhead
Role assignment & Hardware, Software 1, Software 2, Product, Project/Business. Who holds which role. & Now & \\
Directory owners & One owner per top-level directory, matching roles. & After roles & \\
Meeting rhythm & Weekly sync with all five, plus mentor check-ins. How often, which day. & Now & \\
Communication channel & Group chat, issue tracker, or both. & Now & \\
Submission owner & Who finalises the proposal form by 20 Nov. & Week 3 & \\
Mentor expectations (M) & What the mentor reviews, how often, and what they expect to see at each milestone. & Now & \\
\bottomrule
\end{longtable}

\section{Open questions}
\begin{itemize}[nosep]
  \item Do we plan an optional FPGA demo (Hardware role), or keep Demo Day purely simulation-based?
  \item Which users do Product interview first: students, startups, or both?
\end{itemize}

\end{document}
